\documentclass[11pt,twocolumn]{article}
\usepackage[utf8]{inputenc}
\usepackage[T1]{fontenc}
\usepackage[a4paper,margin=2cm]{geometry}
\usepackage{amsmath,amssymb,amsthm,mathtools}
\usepackage[dvipsnames]{xcolor}
\usepackage{booktabs}
\usepackage{hyperref}
\hypersetup{colorlinks=true,linkcolor=blue!70!black,citecolor=green!50!black}
\usepackage{enumitem}
\setlist{nosep,leftmargin=1.5em}
\usepackage{tikz}
\usetikzlibrary{trees,positioning}

\theoremstyle{plain}
\newtheorem{theorem}{Theorem}
\newtheorem{proposition}[theorem]{Proposition}
\newtheorem{corollary}[theorem]{Corollary}
\newtheorem{lemma}[theorem]{Lemma}
\theoremstyle{definition}
\newtheorem{definition}[theorem]{Definition}
\newtheorem{example}[theorem]{Example}
\theoremstyle{remark}
\newtheorem{remark}[theorem]{Remark}

\newcommand{\N}{\mathbb{N}}
\newcommand{\Z}{\mathbb{Z}}
\newcommand{\R}{\mathbb{R}}
\newcommand{\Kk}{\mathbb{K}}

\title{\textbf{The Arithmetic of Conjugated Operations:\\A Number-Theoretic Structure on Twin Fields\\with Applications to Information Theory}}
\author{Jocelyn Nemb\'e\\
\small Roots Insights Laboratory, Singapore\\
\small National Institute of Management Sciences, Libreville\\
\small \texttt{jnembe777@gmail.com, jnembe@root-insights.org}}
\date{}

\begin{document}
\maketitle

% ============================================================================
\begin{abstract}
Given a set $E$ with a binary operation $\odot_0$ and an automorphism $\phi: E \to E$, the \emph{conjugated operations} $x \odot_n y = \phi^{-n}(\phi^n(x) \odot_0 \phi^n(y))$ form a $\Z$-indexed hierarchy governed by the base-change identity $x \odot_{n+k} y = \phi^{-k}(\phi^k(x) \odot_n \phi^k(y))$.
We develop the arithmetic of this hierarchy, showing that the divisibility lattice of $\Z$ induces a lattice structure on operations: GCD gives the \emph{maximal common generator}, LCM gives the \emph{minimal common fusion}, and B\'ezout's identity provides explicit reconstruction.
Beyond translating classical number theory, we prove three \textbf{new results}: (1) a \emph{Composition Theorem} showing that the composition product $\odot_a \circledast \odot_b = \odot_{ab}$ is a genuine monoidal structure on $(E, \odot_0)$ with non-trivial algebraic consequences (Theorem~\ref{thm:composition}); (2) a \emph{differential} entropy--index relation: for a continuous source on $\R$, the differential entropy under $\odot_n$-coding satisfies $h_n = h_0 - n\,I(\phi;X)$ in the $\phi$-stationary regime, where $I(\phi;X) = -\mathbb{E}[\ln|\phi'(X)|]$ is the mean log-contraction of $\phi$ (Theorem~\ref{thm:entropy-index}); in the \emph{discrete} case entropy is bijection-invariant ($H_n \equiv H_0$), so the effect is a genuinely continuous-space phenomenon; (3) a \emph{prime-factorization} of the operation hierarchy: the composition monoid $(\{\odot_n\}_{n\geq1},\circledast)\cong(\N^*,\times)$ is freely generated by the prime operations $\odot_p$, with $\phi^n$ realized by \emph{iterated composition} rather than a Chinese-remainder or tensor-product splitting of fields (Theorem~\ref{thm:prime-decomp}).
We illustrate with three concrete families: dyadic ($\phi(x) = 2x$), exponential ($\phi(x) = e^x$), and relativistic ($\phi(x) = \mathrm{artanh}(x/c)$, yielding Einstein velocity addition).
\end{abstract}

\noindent\textbf{Keywords:} conjugated operations, twin field, divisibility lattice, orbit entropy, prime operations, B\'ezout identity


% ============================================================================
\section{Introduction}
\label{sec:intro}
% ============================================================================

The \emph{twin operations} $\odot_n$ arise naturally in the theory of twin spaces~\cite{nembe2026vol1,nembe2026vol3}: given an automorphism $\phi$ of a structure $(E, \odot_0)$, the conjugated operation $x \odot_n y = \phi^{-n}(\phi^n(x) \odot_0 \phi^n(y))$ transfers the ``level-0'' operation to ``level $n$.'' This construction appears in:

\begin{itemize}
    \item \textbf{Information theory}~\cite{nembe2026vol16}: the group-parameterized entropy $H_G(X)$ depends on which operation $\odot_n$ is used for coding.
    \item \textbf{Relativity}~\cite{nembe2026vol17}: Einstein velocity addition is $\odot_n$ with $\phi = \mathrm{artanh}$.
    \item \textbf{Signal processing}: the $q$-analog operations ($q$-addition, $q$-exponential) are conjugated operations with $\phi(x) = q^x$.
\end{itemize}

The integer $n$ indexes the hierarchy. A natural question: \emph{how does the arithmetic of $\Z$ (divisibility, primes, GCD, congruences) manifest in the algebra of operations?}

We show that the answer goes beyond relabeling. The composition product $\circledast$, the entropy connection, and the prime decomposition are genuinely new algebraic structures that exploit the interplay between $\Z$-arithmetic and the nonlinearity of $\phi$.

Two arithmetics of the index coexist and must not be confused. Composing the gauge maps \emph{adds} indices, $\phi^a\circ\phi^b=\phi^{a+b}$, the additive monoid $(\Z,+)$; the composition product $\circledast$ of \S\ref{sec:composition} instead \emph{multiplies} them, $\odot_a\circledast\odot_b=\odot_{ab}$, the monoid $(\N^*,\times)$. It is the multiplicative structure --- together with the nonlinearity of $\phi$ --- that imports primes, divisibility, and unique factorization into the algebra of operations.


% ============================================================================
\section{The Framework}
\label{sec:framework}
% ============================================================================

\begin{definition}[Twin Hierarchy]\label{def:hierarchy}
Let $(E, \odot_0)$ be a magma (set with binary operation) and $\phi: E \to E$ a bijection. The \textbf{twin hierarchy} is the family $\{\odot_n\}_{n \in \Z}$ defined by:
\begin{equation}\label{eq:twin-op}
\boxed{x \odot_n y = \phi^{-n}(\phi^n(x) \odot_0 \phi^n(y))}
\end{equation}
\end{definition}

\begin{theorem}[Base-Change Identity]\label{thm:base-change}
For all $n, k \in \Z$:
\begin{equation}\label{eq:base-change}
x \odot_{n+k} y = \phi^{-k}(\phi^k(x) \odot_n \phi^k(y))
\end{equation}
The twin hierarchy is governed by a single automorphism: $\odot_n$ is related to $\odot_m$ by conjugation with $\phi^{m-n}$.
\end{theorem}

\begin{proof}
$\phi^k(x) \odot_n \phi^k(y) = \phi^{-n}(\phi^n(\phi^k(x)) \odot_0 \phi^n(\phi^k(y))) = \phi^{-n}(\phi^{n+k}(x) \odot_0 \phi^{n+k}(y))$. Applying $\phi^{-k}$: $\phi^{-k}(\phi^k(x) \odot_n \phi^k(y)) = \phi^{-(n+k)}(\phi^{n+k}(x) \odot_0 \phi^{n+k}(y)) = x \odot_{n+k} y$.
\end{proof}


\subsection{Three Concrete Families}

\begin{example}[Dyadic: $\phi(x) = 2x$, $\odot_0 = +$]\label{ex:dyadic}
$x \odot_n y = 2^{-n}(2^n x + 2^n y) = x + y$. The operations are \emph{all identical}: the dyadic case is \textbf{trivial} because $\phi$ is linear and commutes with $+$. This shows that the arithmetic of twin operations is interesting only when $\phi$ is \textbf{nonlinear}.
\end{example}

\begin{example}[Exponential: $\phi(x) = e^x$, $\odot_0 = +$]\label{ex:exponential}
$x \odot_1 y = \ln(e^x + e^y)$ (the log-sum-exp, fundamental in machine learning).
$x \odot_2 y = \ln(\ln(e^{e^x} + e^{e^y}))$ (double-exponential addition).
Each level $n$ is a more ``nonlinear'' addition. Level 0 is linear; higher levels are increasingly curved.
\end{example}

\begin{example}[Relativistic: $\phi(x) = \mathrm{artanh}(x/c)$, $\odot_0 = +$]\label{ex:relativistic}
$x \odot_1 y = c \tanh(\mathrm{artanh}(x/c) + \mathrm{artanh}(y/c)) = \frac{x + y}{1 + xy/c^2}$.

This is \textbf{Einstein's velocity addition}~\cite{nembe2026vol17}. Higher levels iterate the rapidity map. In natural units ($c = 1$, so $\phi = \mathrm{artanh}$), $\odot_2$ requires $|\mathrm{artanh}(x)| < 1$ and is therefore defined only on the smaller domain $|x| < \tanh(1) \approx 0.76$: each level \emph{contracts} the admissible velocity range rather than the output. (As stated, $\phi$ is not dimensionally homogeneous, so iteration is meaningful only in the dimensionless rapidity formulation.)
\end{example}

\begin{remark}[When is the hierarchy non-trivial?]
The hierarchy is trivial ($\odot_n = \odot_0$ for all $n$) iff $\phi$ is an automorphism of $(E, \odot_0)$. Non-trivial hierarchies require $\phi$ to be a bijection that is \emph{not} a homomorphism --- i.e., $\phi(x \odot_0 y) \neq \phi(x) \odot_0 \phi(y)$.
\end{remark}


% ============================================================================
\section{The Divisibility Lattice of Operations}
\label{sec:lattice}
% ============================================================================

\begin{definition}[Operation Lattice]\label{def:lattice}
The set of operations $\{\odot_n\}_{n \geq 1}$ with:
\begin{align}
\odot_a \wedge \odot_b &:= \odot_{\gcd(a,b)} \quad \text{(maximal common generator)} \\
\odot_a \vee \odot_b &:= \odot_{\mathrm{lcm}(a,b)} \quad \text{(minimal common fusion)}
\end{align}
is a lattice isomorphic to $(\N^*, \mid)$ (the positive integers ordered by divisibility), \emph{provided the hierarchy is faithful}, i.e.\ $n \mapsto \odot_n$ is injective. Otherwise it is the quotient lattice $\N^*/{\sim}$, where $a \sim b$ iff $\odot_a = \odot_b$; the dyadic family (Example~\ref{ex:dyadic}) collapses to a single point.
\end{definition}

\begin{theorem}[Lattice Properties]\label{thm:lattice}
\begin{enumerate}[label=(\alph*)]
    \item \textbf{Generator:} If $d = \gcd(a,b)$, then $\odot_a$ and $\odot_b$ are both expressible as conjugations of $\odot_d$: $\odot_a = \phi^{-d(a/d - 1)} \circ \odot_d \circ (\phi^{d(a/d-1)} \times \phi^{d(a/d-1)})$.
    \item \textbf{Fusion:} If $m = \mathrm{lcm}(a,b)$, then $\odot_m$ is the lowest-index operation expressible as a conjugation of both $\odot_a$ and $\odot_b$.
    \item \textbf{Duality:} $\gcd(a,b) \cdot \mathrm{lcm}(a,b) = ab$, so $(\odot_a \wedge \odot_b) \circledast (\odot_a \vee \odot_b) = \odot_a \circledast \odot_b$ (where $\circledast$ is defined in \S\ref{sec:composition}).
    \item \textbf{B\'ezout reconstruction:} If $au + bv = \gcd(a,b)$, then $\phi^{\gcd(a,b)} = (\phi^a)^u \circ (\phi^b)^v$, allowing $\odot_{\gcd(a,b)}$ to be reconstructed from $\phi^a$ and $\phi^b$.
\end{enumerate}
\end{theorem}

\begin{proof}
(a) follows from the base-change identity with $n = d$, $k = d(p-1)$ where $a = dp$.
(b) follows from $m/a \in \Z$ and $m/b \in \Z$ (by definition of lcm).
(c) is the classical identity $\gcd \cdot \mathrm{lcm} = ab$ applied to indices.
(d) B\'ezout: $\exists u,v \in \Z: au + bv = d$, so $\phi^d = \phi^{au+bv} = (\phi^a)^u \circ (\phi^b)^v$.
\end{proof}


% ============================================================================
\section{The Composition Product (New Result 1)}
\label{sec:composition}
% ============================================================================

\begin{definition}[Composition Product]\label{def:circledast}
The \textbf{composition product} of two operations is:
\begin{equation}
\boxed{\odot_a \circledast \odot_b := \odot_{ab}}
\end{equation}
Explicitly: $\odot_{ab}(x,y) = \phi^{-ab}(\phi^{ab}(x) \odot_0 \phi^{ab}(y))$.

\begin{remark}[$\circledast$ is not function composition]\label{rem:two-monoids}
The product $\circledast$ multiplies levels, whereas composing the underlying maps adds them: $\phi^a\circ\phi^b=\phi^{a+b}$. Thus $(\Z,+)$ (composition of gauge maps) and $(\N^*,\times)$ (the product $\circledast$) are different monoid structures on the same index set, and only the latter carries the prime/divisibility content of \S\ref{sec:prime}. In particular $\odot_a\circledast\odot_b=\odot_{ab}$ is \emph{not} obtained by composing the maps $\phi^a,\phi^b$.
\end{remark}
\end{definition}

\begin{theorem}[Composition is a Monoidal Structure]\label{thm:composition}
The composition product $\circledast$ makes $(\{\odot_n\}_{n \geq 1}, \circledast)$ a commutative monoid with:
\begin{enumerate}[label=(\alph*)]
    \item \textbf{Identity:} $\odot_1$ (since $\odot_1 \circledast \odot_a = \odot_{1 \cdot a} = \odot_a$).
    \item \textbf{Commutativity:} $\odot_a \circledast \odot_b = \odot_{ab} = \odot_{ba} = \odot_b \circledast \odot_a$.
    \item \textbf{Associativity:} $(\odot_a \circledast \odot_b) \circledast \odot_c = \odot_{abc} = \odot_a \circledast (\odot_b \circledast \odot_c)$.
    \item \textbf{Unique factorization:} $\odot_n = \odot_{p_1}^{\circledast \alpha_1} \circledast \cdots \circledast \odot_{p_r}^{\circledast \alpha_r}$ where $n = p_1^{\alpha_1} \cdots p_r^{\alpha_r}$.
\end{enumerate}

As an abstract monoid this is simply $(\N^*,\times)$ pulled back along $n\mapsto\odot_n$; the product $\circledast$ is \emph{defined} by index multiplication, so (a)--(d) are restatements of the arithmetic of $\N^*$. Its concrete content, when $\phi$ is nonlinear, is the functional identity obtained from the base-change theorem (Thm.~\ref{thm:base-change}) with $n=a$, $k=a(b-1)$:
\begin{equation}
(x \odot_{ab} y) = \phi^{-a(b-1)}\!\left(\phi^{a(b-1)}(x) \odot_a \phi^{a(b-1)}(y)\right),
\end{equation}
which expresses level $ab$ as a single conjugation of level $a$ by $\phi^{a(b-1)}$.
\end{theorem}

\begin{proof}
(a)--(c) follow from the multiplicativity of indices.
(d) is the fundamental theorem of arithmetic applied to the index $n$.
The concrete identity $\odot_{ab}(x,y) = \phi^{-a(b-1)}(\phi^{a(b-1)}(x) \odot_a \phi^{a(b-1)}(y))$ on $E$ is precisely the base-change theorem (Thm.~\ref{thm:base-change}) specialized to $n=a$, $k=a(b-1)$.
\end{proof}

\begin{example}[Exponential family]
For $\phi(x) = e^x$: $\odot_6 = \odot_2 \circledast \odot_3$. Explicitly:
\begin{align*}
x \odot_6 y &= \ln^{(6)}(\exp^{(6)}(x) + \exp^{(6)}(y)) \\
&= \phi^{-4}(\phi^4(x) \odot_2 \phi^4(y))
\end{align*}
where $\odot_2(u,v) = \ln(\ln(e^{e^u} + e^{e^v}))$. This is a non-trivial functional identity relating iterated logarithms and exponentials.
\end{example}


% ============================================================================
\section{Connection to Differential Entropy (New Result 2)}
\label{sec:entropy}
% ============================================================================

\begin{theorem}[Differential Entropy--Index Relation]\label{thm:entropy-index}
Let $E = \R$, let $X$ be a \emph{continuous} random variable with differential entropy $h_0 = h(X)$, and let $\phi$ be a $C^1$ diffeomorphism. Under $\odot_n$-coding (codewords formed via $\phi^n$), the effective source is $\phi^n(X)$ and its differential entropy is
\begin{equation}
\boxed{h_n(X) = h(\phi^n(X)) = h_0 - \sum_{k=0}^{n-1} I(\phi; \phi^k(X))}
\end{equation}
where $I(\phi; Y) := h(Y) - h(\phi(Y)) = -\mathbb{E}\!\left[\ln|\phi'(Y)|\right]$ is the \textbf{mean log-contraction} of $\phi$. In the $\phi$-stationary (or ergodic) regime, where $\frac{1}{n}\sum_{k=0}^{n-1} I(\phi;\phi^k(X)) \to I(\phi;X)$, this reduces to the linear law $h_n = h_0 - n\,I(\phi;X)$.
\end{theorem}

\begin{remark}[Why the discrete case is trivial]\label{rem:discrete-trivial}
For a \emph{discrete} $X$, Shannon entropy is invariant under the bijection $\phi$ (a bijection merely relabels outcomes), so $H(\phi^n(X)) = H(X)$ and $H_n \equiv H_0$ for all $n$: no compression occurs. An index-dependent entropy therefore requires the continuous setting above, where the Jacobian $\phi'$ enters through the change-of-variables formula. This is why we work with differential entropy on $\R$.
\end{remark}

\begin{proof}
By the change-of-variables formula for differential entropy, $h(\phi(Y)) = h(Y) + \mathbb{E}[\ln|\phi'(Y)|]$ for any $C^1$ diffeomorphism $\phi$. Telescoping over $k = 0, \dots, n-1$ with $Y_k = \phi^k(X)$,
\[
h(\phi^n X) = h(X) + \sum_{k=0}^{n-1} \mathbb{E}[\ln|\phi'(\phi^k X)|] = h_0 - \sum_{k=0}^{n-1} I(\phi; \phi^k X).
\]
Under stationarity the Birkhoff average $\frac{1}{n}\sum_k I(\phi;\phi^k X) \to I(\phi;X)$, yielding the linear law.
\end{proof}

\begin{example}[Gaussian under a linear contraction]\label{ex:gauss-entropy}
Take $X\sim\mathcal N(0,\sigma^2)$ on $E=\R$ and $\phi(x)=x/2$. Then $\phi'(x)=\tfrac12$, so the per-level contraction is constant, $I(\phi;\phi^kX)=-\mathbb E[\ln\tfrac12]=\ln 2>0$, and the telescoping sum of Theorem~\ref{thm:entropy-index} is exact (no stationarity approximation):
\[
h_n=h_0-n\ln 2,\qquad h_0=\tfrac12\ln(2\pi e\,\sigma^2).
\]
Each level halves the scale and removes $\ln 2$ of differential entropy. Note that $\phi$ here is linear, so the \emph{operations} $\odot_n$ all coincide (Example~\ref{ex:dyadic}), yet the coding source $\phi^n(X)$ --- and hence $h_n$ --- genuinely changes: the entropy law sees $\phi^n$ directly, not the operation it induces. A nonlinear $\phi$ makes both the operations \emph{and} the per-level contraction $I(\phi;\phi^kX)$ vary, and the clean linear law then holds only on the Birkhoff average.
\end{example}

\begin{corollary}[Monotone contraction]
If $I(\phi;X) > 0$ (i.e.\ $\phi$ is, on average, contracting: $\mathbb{E}[\ln|\phi'(X)|] < 0$), then $h_n$ decreases linearly in $n$ and $h_n \to -\infty$: higher-level codes concentrate the source on ever-smaller scales. There is no finite ``critical index''; instead, $n$ trades scale resolution against the description cost of the level (\S\ref{sec:application}).
\end{corollary}

\begin{remark}[Connection to Vol.~16]
When $\phi$ has \emph{finite order} $N$, $G = \langle\phi\rangle \cong \Z/N\Z$ acts on $E$ and the orbit/quotient decomposition $h(X) = h([X]_G) + h(X \mid [X]_G)$ of~\cite{nembe2026vol16} applies, with the index $n \in \Z/N\Z$ parameterizing the depth of symmetry exploitation. For $\phi$ of infinite order ($G \cong \Z$) the present hierarchy is the non-compact analogue: $n$ ranges over $\Z$ and Theorem~\ref{thm:entropy-index} replaces the finite orbit sum by the telescoping log-Jacobian sum.
\end{remark}


% ============================================================================
\section{Prime Factorization of Operations (New Result 3)}
\label{sec:prime}
% ============================================================================

\begin{definition}[Twin Field at Level $n$]\label{def:twin-field}
The \textbf{twin field} $\Kk_n = (E, \odot_n, \otimes_n)$ is the algebraic structure obtained by transferring both addition and multiplication via $\phi^n$:
\begin{align}
x \oplus_n y &= \phi^{-n}(\phi^n(x) + \phi^n(y)) \\
x \otimes_n y &= \phi^{-n}(\phi^n(x) \times \phi^n(y))
\end{align}
When $(E, +, \times) = \Kk_0$ is a field, each $\Kk_n$ is a field isomorphic to $\Kk_0$ via $\phi^n$.
\end{definition}

\begin{definition}[Prime Operation]\label{def:prime-op}
An operation $\odot_p$ is \textbf{prime} if $p$ is a prime number. It cannot be expressed as $\odot_a \circledast \odot_b$ with $a, b > 1$.
\end{definition}

\begin{theorem}[Prime Factorization in the Composition Monoid]\label{thm:prime-decomp}
The composition monoid $(\{\odot_n\}_{n\geq1}, \circledast) \cong (\N^*, \times)$ is the free commutative monoid on the prime operations $\{\odot_p : p \text{ prime}\}$. Hence every $\odot_n$ with $n = p_1^{\alpha_1} \cdots p_r^{\alpha_r}$ factors uniquely as
\begin{equation}
\boxed{\odot_n = \odot_{p_1}^{\circledast \alpha_1} \circledast \cdots \circledast \odot_{p_r}^{\circledast \alpha_r}}
\end{equation}
On the underlying set, this factorization is realized by \emph{iterated composition}: $\phi^{ab} = (\phi^a)^{\circ b}$ (iterate $\phi^a$ a total of $b$ times), and composing these for the prime-power factors reconstructs $\phi^n$.

This is \emph{not} a Chinese-remainder splitting --- the index $n$ is an exponent in $\Z$, not a residue mod $n$ --- and \emph{not} a tensor product of fields: since each $\Kk_n$ is isomorphic to $\Kk_0$ via $\phi^n$ (Definition~\ref{def:twin-field}), all the $\Kk_n$ are mutually isomorphic, and a tensor product of fields is generally not even a field.
\end{theorem}

\begin{proof}
$(\{\odot_n\}, \circledast) \cong (\N^*, \times)$ by Definition~\ref{def:circledast}, and $(\N^*, \times)$ is the free commutative monoid on the primes (fundamental theorem of arithmetic); this gives the unique factorization. The realization $\phi^{ab} = (\phi^a)^{\circ b}$ is the definition of iteration. Crucially, \emph{composition adds exponents}, $\phi^a \circ \phi^b = \phi^{a+b}$, so the prime-power factors combine by \emph{iteration} ($\phi^{ab} = (\phi^a)^{\circ b}$), \emph{not} by composition: $\phi^{p_1^{\alpha_1}} \circ \cdots \circ \phi^{p_r^{\alpha_r}} = \phi^{\sum_i p_i^{\alpha_i}} \neq \phi^{n}$ in general (e.g.\ $\phi^2 \circ \phi^3 = \phi^5 \neq \phi^6$).
\end{proof}

\begin{example}[Factorization of $\odot_{60}$]
$60 = 2^2 \cdot 3 \cdot 5$, so in the composition monoid
\[
\odot_{60} = \odot_2^{\circledast 2} \circledast \odot_3 \circledast \odot_5 = \odot_2 \circledast \odot_2 \circledast \odot_3 \circledast \odot_5.
\]
For the exponential family ($\phi = \exp$), $\odot_{60}$ is realized on $E$ by the $60$-fold iterate $\phi^{60} = \exp^{(60)}$, reached by iterated composition of the prime-power levels (e.g.\ $\phi^4 = (\phi^2)^{\circ 2}$). The tree below is the factorization of the \emph{index} $60$ in $(\N^*, \times)$; its leaves are the prime operations $\odot_2, \odot_3, \odot_5$.

\begin{center}
\begin{tikzpicture}[level distance=1.2cm,
  level 1/.style={sibling distance=3cm},
  level 2/.style={sibling distance=1.5cm},
  every node/.style={font=\small}]
\node {$\odot_{60}$}
  child {node {$\odot_4$}
    child {node {$\odot_2$}}
    child {node {$\odot_2$}}
  }
  child {node {$\odot_3$}}
  child {node {$\odot_5$}};
\end{tikzpicture}
\end{center}
\end{example}


% ============================================================================
\section{Classical Arithmetic in Twin Language}
\label{sec:classical}
% ============================================================================

The classical results of number theory acquire operational meaning:

\begin{proposition}[B\'ezout for Operations]\label{prop:bezout}
If $\gcd(a,b) = d$ with $au + bv = d$, then $\odot_d$ can be reconstructed from $\phi^a$ and $\phi^b$: $\phi^d = (\phi^a)^u \circ (\phi^b)^v$.
\end{proposition}

\begin{proposition}[Coprime Operations]\label{prop:coprime}
$\odot_a$ and $\odot_b$ are \textbf{coprime} ($\gcd(a,b) = 1$) iff the only common generator is $\odot_1$. Coprime operations generate the entire hierarchy: $\forall n \in \Z$, $\odot_n$ is expressible via $\phi^a$ and $\phi^b$.
\end{proposition}

\begin{proposition}[Congruence Classes]\label{prop:congruence}
Suppose $\phi$ has \emph{finite order} $N$ (so $\phi^N = \mathrm{id}$). Then $\odot_m = \odot_{m \bmod N}$, the operation index is well-defined in $\Z/N\Z$, and $a \equiv b \pmod{N}$ implies $\odot_a = \odot_b$. Equivalently, $\odot_b = \phi^{-(b-a)} \circ \odot_a \circ (\phi^{b-a} \times \phi^{b-a})$ depends only on $(b-a) \bmod N$. Without a finite-order hypothesis there is no such periodicity: distinct exponents in $\Z$ generally yield distinct operations.
\end{proposition}

\begin{proposition}[Euler--Fermat Periodicity]\label{prop:euler}
Suppose $\phi$ has finite order $N$, so that $\odot_m = \odot_{m \bmod N}$ and both the additive (composition) and multiplicative ($\circledast$) index arithmetic descend to $\Z/N\Z$. If $\gcd(a, N) = 1$, then by Euler's theorem $a^{\varphi(N)} \equiv 1 \pmod{N}$, hence the $\circledast$-powers of $\odot_a$ are periodic:
\[
\odot_a^{\circledast \varphi(N)} = \odot_{a^{\varphi(N)}} = \odot_1.
\]
Thus $\odot_a$ is invertible in the composition monoid, with order dividing $\varphi(N)$.
\end{proposition}

\begin{remark}[Not tautological]
These propositions \emph{are} translations of classical results. Their value is not novelty but \textbf{interpretation}: B\'ezout reconstruction (Prop.~\ref{prop:bezout}) and coprime generation (Prop.~\ref{prop:coprime}) hold for \emph{any} bijection $\phi$, since they concern the exponents of $\phi$ additively. The congruence and Euler--Fermat statements (Prop.~\ref{prop:congruence}--\ref{prop:euler}), by contrast, require $\phi$ to have \emph{finite order} $N$ --- only then does the index live in $\Z/N\Z$ and acquire genuine modular (periodic) behaviour; for $\phi$ of infinite order they have no operational content, the index being an unbounded exponent in $\Z$.
\end{remark}


% ============================================================================
\section{Application: Optimal Operation Index for Compression}
\label{sec:application}
% ============================================================================

\begin{theorem}[Optimal Coding Level]\label{thm:optimal-level}
Let $X$ be a continuous source on $\R$ with differential entropy $h_0$, encoded at a fixed target distortion $\delta$, and let $\{\odot_n\}$ have $I(\phi;X) > 0$. Write $R_n(\delta) = h_n(X) + \log_2(1/\delta_n) + C_n$ for the rate at level $n$, where $\delta_n$ is the level-$n$ cell size and $C_n$ is the description cost of the index $\odot_n$.
\begin{enumerate}[label=(\alph*)]
    \item The optimal coding operation is $\odot_{n^*}$ with $n^* = \arg\min_n R_n(\delta)$.
    \item For prime $n = p$: $\odot_p$ cannot be decomposed, so $C_p$ pays the full index cost. For composite $n = \prod p_i^{\alpha_i}$: $\odot_n$ is described by its prime factorization (Theorem~\ref{thm:prime-decomp}), so $C_n = \sum_i \alpha_i \log_2 p_i$.
    \item The trade-off: higher $n$ lowers $h_n$ but, under $\phi^n$, also rescales the cells $\delta_n$; the rate $h_n + \log_2(1/\delta_n)$ is scale-invariant, so the genuine optimization is over the index cost $C_n$ at the target distortion.
\end{enumerate}
\end{theorem}

\begin{remark}[The unconstrained problem is degenerate]
Minimizing $h_n + C_n$ \emph{without} fixing a distortion has no finite optimum when $I(\phi;X) > 0$: by the Monotone-contraction corollary $h_n \to -\infty$ linearly while $C_n$ grows only logarithmically, so $n \to \infty$ formally. This is an artifact of differential entropy under coordinate contraction; fixing the distortion $\delta$ (equivalently, tracking the cell size $\delta_n$) restores scale-invariance and a well-posed optimum, as in~(c).
\end{remark}

\begin{remark}[Connection to optimal group selection]
This is the twin-operation analog of the optimal group selection problem in~\cite{nembe2026a2}: choosing the ``right'' operation level $n$ (equivalently, the ``right'' group $\langle\phi^n\rangle \leq \langle\phi\rangle$) maximizes an information-theoretic criterion.
\end{remark}


% ============================================================================
\section{Conclusion}
\label{sec:conclusion}
% ============================================================================

We have developed the arithmetic of conjugated operations, going beyond the translation of classical number theory into twin notation. The three new results are:

\begin{enumerate}
    \item The \textbf{composition product} $\circledast$ is a monoidal structure with non-trivial functional-equation content when $\phi$ is nonlinear (Thm.~\ref{thm:composition}).
    \item The \textbf{differential entropy--index relation} $h_n = h_0 - n\,I(\phi;X)$ (in the $\phi$-stationary regime, with $I = -\mathbb{E}[\ln|\phi'|]$) links the operation level to scale contraction; in the discrete case Shannon entropy is bijection-invariant, so the effect is purely a continuous-space phenomenon (Thm.~\ref{thm:entropy-index}).
    \item The \textbf{prime factorization} of the operation hierarchy: the composition monoid $(\{\odot_n\}, \circledast) \cong (\N^*, \times)$ is free on the prime operations, with $\phi^n$ realized by iterated composition --- not by a Chinese-remainder or tensor-product splitting of fields (Thm.~\ref{thm:prime-decomp}).
\end{enumerate}

These results, combined with the classical arithmetic (B\'ezout, Euler, lattice), provide a complete number-theoretic framework for the study of twin operations, with applications to coding theory, signal processing, and mathematical physics.

\textbf{Open problems:} (1) Extension to non-cyclic groups $G$ (non-monogenic hierarchies). (2) The analog of the prime number theorem for the density of prime operations in continuous hierarchies. (3) The finite-order case ($\phi^N = \mathrm{id}$), where the index lives in $\Z/N\Z$ and genuine congruence and Galois-type phenomena on the composition monoid $(\Z/N\Z, \times)$ appear (cf.\ Propositions~\ref{prop:congruence}--\ref{prop:euler}).


% ============================================================================
\begin{thebibliography}{10}
\bibitem{nembe2026vol1} J.~Nemb\'e, ``Twin Spaces: Foundations,'' \emph{ROOTS-INSIGHTS}, Vol.~1, 2026.
\bibitem{nembe2026vol3} J.~Nemb\'e, ``The Transfer Theorem,'' \emph{ROOTS-INSIGHTS}, Vol.~3, 2026.
\bibitem{nembe2026vol16} J.~Nemb\'e, ``Twin Information Theory,'' \emph{ROOTS-INSIGHTS}, Vol.~16, 2026.
\bibitem{nembe2026vol17} J.~Nemb\'e, ``Twin General Relativity,'' \emph{ROOTS-INSIGHTS}, Vol.~17, 2026.
\bibitem{nembe2026a2} J.~Nemb\'e, ``Optimal Group Selection for Communication Systems,'' 2026.
\bibitem{hardy1979} G.H.~Hardy, E.M.~Wright, \emph{An Introduction to the Theory of Numbers}, 5th ed., Oxford, 1979.
\bibitem{kaneko2004} M.~Kaneko, ``Poly-Bernoulli numbers,'' \emph{J.\ Th\'eorie Nombres Bordeaux}, vol.~9, pp.~221--228, 1997.
\bibitem{borwein2013} P.~Borwein, S.~Choi, B.~Rooney, A.~Weirathmueller, \emph{The Riemann Hypothesis}, Springer, 2008.
\end{thebibliography}

\end{document}
